\chapter{Monotone maps and singular inverse derivatives}\label{sec:monotone-inverse-derivatives}
The area formula can detect a singular part of an inverse even when the original map is Lipschitz. We derive this one-dimensional application from the area formula of \cite{Sim18}, Chapter 2, Section 3, and the interval BV representation already proved in Section \ref{sec:bv-interval-structure}. Throughout $a<b$, $f:[a,b]\to\mathbb R$ is strictly increasing and Lipschitz, $c=f(a)$, $d=f(b)$, and $g:[c,d]\to[a,b]$ is its continuous inverse. Endpoint values do not carry Lebesgue mass.

\begin{lemma}[The inverse derivative is a pushforward of interval length]\label{lem:monotone-inverse-pushforward}
The positive measure $\mu=f_\#(\mathcal L^1\llcorner[a,b])$ has no atoms, is supported on $[c,d]$, and has mass $b-a$. For $c\leq u<v\leq d$,
\[
 \mu((u,v])=g(v)-g(u),
 \qquad Dg=\mu\llcorner(c,d).
\]
The distributional derivative is taken on the open interval $(c,d)$.
\end{lemma}
\begin{proof}
Strict increase and continuity give $f^{-1}((u,v])=(g(u),g(v)]$, so pushing forward length gives the increment identity. Every target singleton has at most one preimage and therefore zero mass. Support and total mass follow directly from the definition of pushforward. In particular $g(t)=a+\mu((c,t])$ for $c\leq t\leq d$. Lemma \ref{lem:interval-measure-primitive} identifies its distributional derivative on $(c,d)$ with $\mu$. The endpoint atoms vanish, so the total mass of this derivative is also $b-a$. This establishes the BV assertion directly, without assuming absolute continuity of the inverse.
\end{proof}

\begin{theorem}[Regular and singular parts of a monotone inverse]\label{thm:monotone-inverse-derivative-decomposition}
Let $R\subset(a,b)$ be the differentiability points of $f$ where $f'>0$, and put $C=(a,b)\setminus R$. Define
\[
 h(y)=\frac1{f'(g(y))}\quad\text{for }y\in f(R),
 \qquad h(y)=0\quad\text{otherwise}.
\]
Then $h$ is measurable and integrable, and
\[
 Dg=h\mathcal L^1+\nu,
 \qquad \nu=f_\#(\mathcal L^1\llcorner C),
 \qquad \nu\perp\mathcal L^1.
\]
Here the measures are restricted to $(c,d)$. Their masses are $\int_c^dh=|R|$ and $\nu((c,d))=|C|$. The singular measure $\nu$ has no atoms.
\end{theorem}
\begin{proof}
A Lipschitz function is differentiable almost everywhere, and monotonicity gives a nonnegative derivative wherever it exists. The differentiability set and derivative admit Borel representatives, by expressing differentiability through countably many rational difference quotients. Use these representatives throughout; the continuous strictly monotone map is a homeomorphism of its compact interval onto its image, so it also carries Borel sets to Borel sets. Extend $f$ constantly outside $[a,b]$ to apply the whole-line area formula; this preserves its Lipschitz bound and changes no interior derivative. Apply the weighted area formula, Proposition \ref{prop:weighted-area}, on $R$ with weight $\chi_{f^{-1}(B)}/f'$ for a Borel target set $B$. The weight is finite at each point of $R$, although not necessarily bounded. Injectivity makes the counting sum consist of the single inverse point when it lies in $R$. The resulting equality is
\[
 |R\cap f^{-1}(B)|=\int_Bh(y)\,dy.
\]
Thus the regular pushforward has density $h$, with total mass $|R|$; measurability follows from the weighted formula or from the measurable derivative composed with the continuous inverse.

On the differentiability points of $C$, the derivative vanishes. The remaining points form a Lebesgue-null set. Corollary \ref{cor:area-null-image-criterion} and the Lipschitz null-image bound give $|f(C)|=0$. The critical pushforward is therefore concentrated on a Lebesgue-null set, giving the singular measure $\nu$. Its total mass is $|C|$, and it has no atoms because the restriction of an injective map still has at most one point in any singleton preimage. Finally split the measure in Lemma \ref{lem:monotone-inverse-pushforward} into the restrictions to $R$ and $C$; omitted endpoints have zero length. This proves the decomposition. Singularity here uses the one-dimensional area formula and injectivity; it is not an assertion about critical pushforwards for arbitrary maps between different dimensions.
\end{proof}

\begin{corollary}[An exact criterion for absolute continuity of the inverse]\label{cor:monotone-inverse-ac-criterion}
The inverse $g$ is absolutely continuous on $[c,d]$ if and only if $f'>0$ almost everywhere on $(a,b)$. If this holds, then
\[
 g(y)=a+\int_c^y\frac1{f'(g(s))}\,ds,
 \qquad \int_c^d\frac1{f'(g(s))}\,ds=b-a,
\]
where the reciprocal is assigned any value on its null exceptional set.
\end{corollary}
\begin{proof}
By Proposition \ref{prop:interval-bv-ac-lipschitz}, absolute continuity of the regulated representative of $g$ is equivalent to absolute continuity of $Dg$ with respect to Lebesgue measure. In Theorem \ref{thm:monotone-inverse-derivative-decomposition}, the positive singular part has mass $|C|$. It vanishes precisely when $|C|=0$, or equivalently $f'>0$ almost everywhere. The formula and integral identity then follow from Lemma \ref{lem:interval-measure-primitive} and the total derivative mass $b-a$. Strict monotonicity alone does not ensure this positivity condition.
\end{proof}

\begin{corollary}[A lower derivative bound is exactly an inverse Lipschitz bound]\label{cor:monotone-inverse-lipschitz-criterion}
For $K>0$, the inverse has Lipschitz constant at most $K$ if and only if $f'\geq1/K$ almost everywhere on $(a,b)$.
\end{corollary}
\begin{proof}
If $g$ is $K$-Lipschitz, then for $x<x'$ one has $x'-x\leq K(f(x')-f(x))$. Divide by $x'-x$ and pass to the derivative at almost every point to obtain $f'\geq1/K$. Conversely, Lipschitz regularity of $f$ gives absolute continuity and $f(x')-f(x)=\int_x^{x'}f'$. The assumed lower bound yields $f(x')-f(x)\geq(x'-x)/K$. Put $x=g(y)$ and $x'=g(y')$ to obtain the inverse Lipschitz estimate, including endpoint values by continuity. A positive derivative almost everywhere, without a uniform positive lower bound, gives only the absolute-continuity conclusion of Corollary \ref{cor:monotone-inverse-ac-criterion}.
\end{proof}

\begin{example}[An inverse has derivative one almost everywhere and a singular remainder]\label{ex:fat-cantor-inverse-singular-part}
Let $f:[0,1]\to[0,1/2]$ and $\nu$ be as in Example \ref{ex:area-fat-cantor-critical-map}, and let $g=f^{-1}$. Then
\[
 g(y)=y+\nu([0,y]),\qquad Dg=\mathcal L^1+\nu
 \quad\text{on }(0,1/2).
\]
The function $g$ is continuous and increasing, $g'=1$ almost everywhere, and $g(1/2)-g(0)=1$. It is not absolutely continuous. The function $g(y)-y$ is continuous and nondecreasing, has derivative zero almost everywhere, and increases from zero to $1/2$.
\end{example}
\begin{proof}
Example \ref{ex:area-fat-cantor-critical-map} gives $f_\#\mathcal L^1=\mathcal L^1\llcorner[0,1/2]+\nu$ and $\nu([0,1/2])=1/2$, with no atoms. Lemma \ref{lem:monotone-inverse-pushforward} gives the displayed primitive and derivative. The compact set $f(E)$ has measure zero. On every complementary interval, the inverse lies in an interval disjoint from $E$ where $f'=1$ and $f$ is affine; hence $g$ is affine with slope one there. Thus $g'=1$ off $f(E)$ and in particular almost everywhere. Its derivative integral over $(0,1/2)$ is only $1/2$, whereas its total increment is one. The remaining half is the positive singular measure, so Proposition \ref{prop:interval-bv-ac-lipschitz} rules out absolute continuity. The final assertions follow from $g-y=\nu([0,y])$ and absence of atoms.
\end{proof}

\begin{proposition}[Densities can concentrate into this singular continuous measure]\label{prop:fat-cantor-regularized-critical-density}
For the same $f,E$, let $f_\varepsilon(x)=f(x)+\varepsilon x$, $0<\varepsilon\leq1$. This map has inverse Lipschitz constant at most $1/\varepsilon$. Its restricted critical-mass pushforward satisfies
\[
 \nu_\varepsilon=(f_\varepsilon)_\#(\mathcal L^1\llcorner E)
 =\frac1\varepsilon\chi_{f_\varepsilon(E)}\mathcal L^1,
 \qquad |f_\varepsilon(E)|=\frac\varepsilon2.
\]
Each $\nu_\varepsilon$ is absolutely continuous with mass $1/2$, but $\nu_\varepsilon$ converges weakly to the singular continuous measure $\nu$ as $\varepsilon\downarrow0$.
\end{proposition}
\begin{proof}
The derivative is at least $\varepsilon$ almost everywhere, so Corollary \ref{cor:monotone-inverse-lipschitz-criterion} gives the inverse bound. On $E$, the derivative equals $\varepsilon$ almost everywhere. The injective area formula gives $|f_\varepsilon(E)|=\int_Ef_\varepsilon'=\varepsilon|E|=\varepsilon/2$. Its weighted version with weight $\chi_E/\varepsilon$ gives the density formula, first against Borel indicators and then against arbitrary nonnegative tests.

For a continuous compactly supported test function $\phi$,
\[
 \int\phi\,d\nu_\varepsilon=\int_E\phi(f(x)+\varepsilon x)\,dx
 \longrightarrow\int_E\phi(f(x))\,dx=\int\phi\,d\nu.
\]
Uniform continuity on a common compact interval and $|\varepsilon x|\leq\varepsilon$ justify convergence. The masses remain $1/2$ while the lengths of the density supports tend to zero. The limiting measure has no atoms by Example \ref{ex:area-fat-cantor-critical-map}, so concentration need not produce a Dirac mass.
\end{proof}
